% Bibliografía del capítulo de revisión de la literatura (prefijo literatura-).
% Todas las entradas se verificaron con búsquedas web el 5 de octubre de 2026 (título,
% autores, año y medio). Volumen, páginas y DOI se incluyen solo cuando aparecieron en la
% ficha del editor o de un índice bibliográfico.

% --- (1) CanSat con autogiro
\bibitem{literatura-rivadeneira2023} F.~Rivadeneira, D.~Godinez, K.~Kiyan, V.~Huayapa,
  S.~Acosta, N.~Pérez, A.~Hinostroza y D.~Arce. Run2Sat I: Design and implementation of a
  CanSat with Autogyro System. En \emph{1st IEEE Colombian Caribbean Conference (C3)},
  Barranquilla, 2023.

\bibitem{literatura-rivadeneira2025} F.~Rivadeneira, K.~Kiyan, V.~Huayapa, D.~Godinez,
  N.~Pérez, A.~Hinostroza, S.~Acosta y D.~Arce. To eggspace and beyond: design and
  implementation of an autogyro-CanSat with IoT purposes using AWS. \emph{Transactions on
  Energy Systems and Engineering Applications}, 6(2), 2025.
  doi:10.32397/tesea.vol6.n2.628.

\bibitem{literatura-villacorta2025} L.~Villacorta, F.~Chunga, G.~Aquino y C.~D.~Salvador.
  Enhancing CanSat Mission Safety: An Autogyro-Based Landing Solution. En
  \emph{Proceedings of the 10th Brazilian Technology Symposium (BTSym'24)}, Smart
  Innovation, Systems and Technologies 443, pp.~107--115. Springer, 2025.
  doi:10.1007/978-3-031-92651-8\_12.

\bibitem{literatura-shukla2022} P.~Shukla, R.~Mishra, U.~A.~Sardar y B.~Mohapatra.
  Satellite Design for CANSAT with Autorotatig payloads. En \emph{2022 4th International
  Conference on Advances in Computing, Communication Control and Networking (ICAC3N)},
  pp.~2385--2392. IEEE, 2022. doi:10.1109/ICAC3N56670.2022.10074552.

\bibitem{literatura-kizilkaya2017} M.~Ö.~Kizilkaya, A.~E.~Oğuz y S.~Soyer. CanSat descent
  control system design and implementation. En \emph{8th International Conference on
  Recent Advances in Space Technologies (RAST)}. IEEE, 2017.

\bibitem{literatura-swifturk2019} Equipo SwifTurk (\#1387). \emph{CanSat 2019 Critical
  Design Review (CDR)}. CanSat Competition, 2019. Documento de equipo difundido en
  ResearchGate.

% --- (2) Autorrotación vertical de rotores pequeños y desaceleradores rotativos simples
\bibitem{literatura-lugt1983} H.~J.~Lugt. Autorotation. \emph{Annual Review of Fluid
  Mechanics}, 15:123--147, 1983. doi:10.1146/annurev.fl.15.010183.001011.

\bibitem{literatura-iversen1979} J.~D.~Iversen. Autorotating flat-plate wings: the effect
  of the moment of inertia, geometry and Reynolds number. \emph{Journal of Fluid
  Mechanics}, 92(2):327--348, 1979. doi:10.1017/S0022112079000641.

\bibitem{literatura-seter2014} D.~Seter y A.~Rosen. Theoretical and Experimental Study of
  Axial Autorotation of Simple Rotary Decelerators. \emph{Journal of Aircraft},
  51(1):236--248, 2014. doi:10.2514/1.C032305.

\bibitem{literatura-brindejonc2007} A.~Brindejonc, J.~Sirohi e I.~Chopra. Design and
  Testing of an Autorotative Payload Delivery System. \emph{Journal of the American
  Helicopter Society}, 52(4):360--370, 2007. doi:10.4050/JAHS.52.360.

\bibitem{literatura-jain2022} P.~Jain, J.~Sirohi y C.~G.~Cameron. Design, Analysis, and
  Testing of a Passively Deployable Autorotative Decelerator. \emph{Journal of Aircraft},
  59(1), 2022 (en línea, noviembre de 2021). doi:10.2514/1.C036509.

\bibitem{literatura-nadal2005} V.~Nadal Mora, Á.~Sanz-Andrés y A.~Cuerva. Experimental
  Investigation of an Autorotating-Wing Aerodynamic Decelerator System. \emph{18th AIAA
  Aerodynamic Decelerator Systems Technology Conference and Seminar}, AIAA Paper
  2005-1635, 2005.

\bibitem{literatura-nadal2006} V.~Nadal Mora, Á.~Sanz-Andrés y A.~Cuerva Tejero. Model of
  the Aerodynamic Behavior of a Pararotor. \emph{Journal of Aircraft}, 2006.
  doi:10.2514/1.21435.

\bibitem{literatura-piechocki2014} J.~Piechocki, V.~Nadal Mora y Á.~Sanz-Andrés. Pararotor
  Dynamics: Center of Mass Displacement from the Blade Plane---Analytical Approach.
  \emph{Journal of Aircraft}, 51(2):651--660, 2014. doi:10.2514/1.C032378.

\bibitem{literatura-martiarena2015} J.~F.~Martiarena, V.~Nadal Mora y J.~Piechocki.
  Experimental study of the effect of blade curvature and aspect ratio on the performance
  of a rotary-wing decelerator. \emph{Aerospace Science and Technology}, 43:471--477,
  2015.

\bibitem{literatura-veismann2022} M.~Veismann, D.~Yos y M.~Gharib. Parametric study of
  small-scale rotors in axial descent. \emph{Physics of Fluids}, 34(3):035124, 2022.

% --- (3) Sámaras y robots inspirados en sámaras
\bibitem{literatura-ulrich2010a} E.~R.~Ulrich, D.~J.~Pines y J.~S.~Humbert. From falling
  to flying: the path to powered flight of a robotic samara nano air vehicle.
  \emph{Bioinspiration \& Biomimetics}, 5(4):045009, 2010.
  doi:10.1088/1748-3182/5/4/045009.

\bibitem{literatura-ulrich2010b} E.~R.~Ulrich, J.~S.~Humbert y D.~J.~Pines. Pitch and
  Heave Control of Robotic Samara Micro Air Vehicles. \emph{Journal of Aircraft},
  47:1290--1299, 2010. doi:10.2514/1.47197.

\bibitem{literatura-kellas2007} A.~Kellas. \emph{The Guided Samara: Design and Development
  of a Controllable Single-Bladed Autorotating Vehicle}. Tesis de maestría, Massachusetts
  Institute of Technology, 2007.

\bibitem{literatura-matic2015} G.~Matič, M.~Topič y M.~Jankovec. Mathematical Model of a
  Monocopter Based on Unsteady Blade-Element Momentum Theory. \emph{Journal of Aircraft},
  52(6):1905--1913, 2015. doi:10.2514/1.C033098.

\bibitem{literatura-bai2019} S.~Bai y P.~Chirarattananon. Design and Take-Off Flight of a
  Samara-Inspired Revolving-Wing Robot. En \emph{2019 IEEE/RSJ International Conference on
  Intelligent Robots and Systems (IROS)}, Macao, 2019.

\bibitem{literatura-bai2022} S.~Bai, Q.~He y P.~Chirarattananon. A bioinspired
  revolving-wing drone with passive attitude stability and efficient hovering flight.
  \emph{Science Robotics}, 7(66), 2022. doi:10.1126/scirobotics.abg5913.

\bibitem{literatura-varshney2012} K.~Varshney, S.~Chang y Z.~J.~Wang. The kinematics of
  falling maple seeds and the initial transition to a helical motion.
  \emph{Nonlinearity}, 25:C1--C8, 2012.

\bibitem{literatura-salcedo2013} E.~Salcedo, C.~Treviño, R.~O.~Vargas y
  L.~Martínez-Suástegui. Stereoscopic particle image velocimetry measurements of the
  three-dimensional flow field of a descending autorotating mahogany seed (\emph{Swietenia
  macrophylla}). \emph{Journal of Experimental Biology}, 216(11):2017--2030, 2013.

\bibitem{literatura-lee2014} S.~J.~Lee, E.~J.~Lee y M.~H.~Sohn. Mechanism of autorotation
  flight of maple samaras (\emph{Acer palmatum}). \emph{Experiments in Fluids}, 55:1718,
  2014.

\bibitem{literatura-jung2016} B.~K.~Jung y D.~Rezgui. Investigating the Autorotational
  Performance of Scaled Samara Rotor in Vertical and Forward Flight. En \emph{42nd European
  Rotorcraft Forum}, Lille, 2016.

\bibitem{literatura-jung2024} B.~K.~Jung y D.~Rezgui. Experimental study into the effects
  of pitch and coning angles on the flight performance of the natural samara. \emph{Royal
  Society Open Science}, 11(9):240758, 2024. doi:10.1098/rsos.240758.

\bibitem{literatura-schaeffer2024} B.~M.~Schaeffer, S.~S.~Truman, T.~T.~Truscott y
  A.~K.~Dickerson. Maple samara flight is robust to morphological perturbation and united
  by a classic drag model. \emph{Communications Biology}, 7:248, 2024.
  doi:10.1038/s42003-024-05913-3.

\bibitem{literatura-elmakdah2022} A.~M.~El Makdah, K.~Zhang y D.~E.~Rival. On the robust
  autorotation of a samara-inspired rotor in gusty environments. \emph{Bioinspiration \&
  Biomimetics}, 17(4), 2022. doi:10.1088/1748-3190/ac68bb.

\bibitem{literatura-lolli2026} A.~Lolli, G.~Corsi, B.~Mazzolai y A.~DeSimone.
  Aerodynamic performance of autorotating seeds: scaling by size. \emph{Bioinspiration \&
  Biomimetics}, 21(2):026017, 2026. doi:10.1088/1748-3190/ae4f46.

% --- (4) Bajo número de Reynolds
\bibitem{literatura-laitone1997} E.~V.~Laitone. Wind tunnel tests of wings at Reynolds
  numbers below 70\,000. \emph{Experiments in Fluids}, 23:405--409, 1997.
  doi:10.1007/s003480050128.

\bibitem{literatura-okamoto2011} M.~Okamoto y A.~Azuma. Aerodynamic Characteristics at Low
  Reynolds Numbers for Wings of Various Planforms. \emph{AIAA Journal},
  49(6):1135--1150, 2011. doi:10.2514/1.J050071.

\bibitem{literatura-bruining1979} A.~Bruining. \emph{Aerodynamic characteristics of a
  curved plate airfoil section at Reynolds numbers 60000 and 100000 and angles of attack
  from $-10$ to $+90$ degrees}. Informe LR-281, Delft University of Technology, 1979.

\bibitem{literatura-shyy2008} W.~Shyy, Y.~Lian, J.~Tang, D.~Viieru y H.~Liu.
  \emph{Aerodynamics of Low Reynolds Number Flyers}. Cambridge Aerospace Series,
  Cambridge University Press, 2008. ISBN 978-0-521-88278-1.

% --- (5) Desaceleradores rotativos para entrada y recuperación
\bibitem{literatura-barzda1966} J.~J.~Barzda. Rotors for recovery. \emph{Journal of
  Spacecraft and Rockets}, 3(1):104--109, 1966.

\bibitem{literatura-young2004} L.~A.~Young, G.~A.~Briggs, E.~W.~Aiken y G.~Pisanich.
  Rotary-Wing Decelerators for Probe Descent Through the Atmosphere of Venus. En
  \emph{2nd International Planetary Probe Workshop}, NASA Ames Research Center, 2004.

\bibitem{literatura-diazsilva2013} R.~A.~Diaz-Silva, D.~Arellano, M.~M.~Sarigul-Klijn y
  N.~Sarigul-Klijn. Rotary Decelerators for Spacecraft: Historical Review and Simulation
  Results. \emph{AIAA SPACE 2013 Conference and Exposition}, AIAA Paper 2013-5361, 2013.

\bibitem{literatura-bergmann2024} P.~Bergmann, C.~Riegler, Z.~Klaschka, T.~Herbst,
  J.~M.~Wolf, M.~Reigl y otros. Daedalus 2: Autorotation Entry, Descent and Landing
  Experiment on REXUS29. Preimpresión arXiv:2406.11292, 2024.

\bibitem{literatura-marques2024} J.~P.~C.~S.~Marques, A.~de Souza y F.~S.~Cunha. Design of
  a Helicopter System Recovery for a Sounding Rocket. \emph{International Journal of
  Aeronautical and Space Sciences}, 25:662--675, 2024. doi:10.1007/s42405-023-00696-z.

\bibitem{literatura-patnala2024} S.~Patnala, P.~Elango, R.~Mohan y Shamrao.
  Autorotation-Based Descent with Trajectory Optimization. \emph{Journal of Aircraft},
  61(1), 2024. doi:10.2514/1.C036912.

% --- (6) Anillo de vórtices y estela turbulenta
\bibitem{literatura-glauert1926b} H.~Glauert. The Analysis of Experimental Results in the
  Windmill Brake and Vortex Ring States of an Airscrew. \emph{Aeronautical Research
  Committee, Reports and Memoranda} No.~1026, 1926.

\bibitem{literatura-drees1951} J.~M.~Drees y W.~P.~Hendal. Airflow patterns in the
  neighbourhood of helicopter rotors. \emph{Aircraft Engineering}, 23, 1951.

\bibitem{literatura-washizu1966} K.~Washizu, A.~Azuma, J.~Koo y T.~Oka. Experiments on a
  model helicopter rotor operating in the vortex ring state. \emph{Journal of Aircraft},
  3:225--230, 1966.

\bibitem{literatura-green2005} R.~B.~Green, E.~A.~Gillies y R.~E.~Brown. The flow field
  around a rotor in axial descent. \emph{Journal of Fluid Mechanics}, 534:237--261, 2005.
  doi:10.1017/S0022112005004155.

\bibitem{literatura-stack2005} J.~Stack, F.~X.~Caradonna y Ö.~Savaş. Flow visualizations
  and extended thrust time histories of rotor vortex wakes in descent. \emph{Journal of
  the American Helicopter Society}, 50(3):279--288, 2005.

\bibitem{literatura-shetty2011} O.~R.~Shetty y M.~S.~Selig. Small-Scale Propellers
  Operating in the Vortex Ring State. \emph{49th AIAA Aerospace Sciences Meeting}, AIAA
  Paper 2011-1254, Orlando, 2011. doi:10.2514/6.2011-1254.

% --- (7) BEMT y correcciones de alta inducción
\bibitem{literatura-wilson1974} R.~E.~Wilson y P.~B.~S.~Lissaman. \emph{Applied
  Aerodynamics of Wind Power Machines}. Informe NSF-RA-N-74-113, Oregon State University,
  1974.

\bibitem{literatura-buhl2005} M.~L.~Buhl Jr. \emph{A New Empirical Relationship between
  Thrust Coefficient and Induction Factor for the Turbulent Windmill State}. NREL/TP-500-36834,
  National Renewable Energy Laboratory, 2005.

\bibitem{literatura-moriarty2005} P.~J.~Moriarty y A.~C.~Hansen. \emph{AeroDyn Theory
  Manual}. NREL/TP-500-36881, National Renewable Energy Laboratory, 2005.
  doi:10.2172/15014831.

\bibitem{literatura-shen2005} W.~Z.~Shen, R.~Mikkelsen, J.~N.~Sørensen y C.~Bak. Tip loss
  corrections for wind turbine computations. \emph{Wind Energy}, 8(4):457--475, 2005.
  doi:10.1002/we.153.

\bibitem{literatura-ning2014} S.~A.~Ning. A simple solution method for the blade element
  momentum equations with guaranteed convergence. \emph{Wind Energy}, 17(9):1327--1345,
  2014. doi:10.1002/we.1636.

\bibitem{literatura-hansen2015} M.~O.~L.~Hansen. \emph{Aerodynamics of Wind Turbines},
  3.ª ed. Routledge, 2015. ISBN 978-1-138-77507-7.

\bibitem{literatura-liew2024} J.~Liew, K.~S.~Heck y M.~F.~Howland. Unified momentum model
  for rotor aerodynamics across operating regimes. \emph{Nature Communications}, 15:6658,
  2024. doi:10.1038/s41467-024-50756-5.

% --- (8) Paracaídas pequeños
\bibitem{literatura-ewing1978} E.~G.~Ewing, H.~W.~Bixby y T.~W.~Knacke. \emph{Recovery
  Systems Design Guide}. AFFDL-TR-78-151, Air Force Flight Dynamics Laboratory, 1978.

\bibitem{literatura-cockrell1987} D.~J.~Cockrell. \emph{The Aerodynamics of Parachutes}.
  AGARDograph 295, AGARD/OTAN, 1987.

\bibitem{literatura-scher1971} S.~H.~Scher e I.~G.~Young. \emph{Drag coefficients for
  partially inflated flat circular parachutes}. NASA TN D-6423, 1971.

\bibitem{literatura-johari2005} H.~Johari y K.~J.~Desabrais. Vortex shedding in the near
  wake of a parachute canopy. \emph{Journal of Fluid Mechanics}, 536, 2005.

\bibitem{literatura-mattei2023} M.~Mattei, C.~Phillippe, F.~Panerai y L.~Villafañe-Roca.
  Subscale Wind Tunnel Testing of Parachutes: Effects of Design Shape on Drag and Canopy
  Breathing. \emph{AIAA SCITECH 2023 Forum}, 2023. doi:10.2514/6.2023-0462.

\bibitem{literatura-guglieri2012} G.~Guglieri. Parachute-Payload System Flight Dynamics
  and Trajectory Simulation. \emph{International Journal of Aerospace Engineering},
  2012:182907, 2012. doi:10.1155/2012/182907.
